% =========================================================================
% Capítulo 0: Repaso
% Tema 01: Álgebra de vectores
% Problema 01: Operaciones con Base Única
% =========================================================================
\subsubsection{Operaciones con Base Única}
Se define una base orto-normal $\{\vec{e}_1, \vec{e}_2, \vec{e}_3\}$. Considere $a, b, c$ como escalares constantes. Resuelva:

\begin{multicols}{2}
\begin{enumerate}
    \item $\vec{e}_1 \cdot \vec{e}_1 =$
    \item $a\vec{e}_1 \cdot b\vec{e}_2 =$
    \item $\vec{e}_2 \times \vec{e}_3 =$
    \item $\vec{e}_3 \times \vec{e}_2 =$
    \item $\vec{e}_1 \times \vec{e}_1 =$
    \item $\| a\vec{e}_2 + b\vec{e}_2 \| =$
    \item $(a\vec{e}_1 + b\vec{e}_2) \cdot \vec{e}_1 =$
    \item $\| a\vec{e}_1 \times b\vec{e}_3 \| =$
    \item $\vec{e}_1 \cdot (\vec{e}_2 \times \vec{e}_3) =$
    \item $(\vec{e}_1 \times \vec{e}_2) \cdot \vec{e}_3 =$
    \item $\| a\vec{e}_1 - a\vec{e}_1 \| =$
    \item $(a\vec{e}_1 + b\vec{e}_3) \cdot (c\vec{e}_2) =$
\end{enumerate}
\end{multicols}
\textbf{Segunda parte.}
  
Se define un marco de referencia y a el se fija la siguiente base $\{\Vec{e}_1,\Vec{e}_2,\Vec{e}_3\}$ que es orto-normal, resuelva las siguientes operaciones:
\begin{multicols}{2}
\begin{enumerate}
    \item$\Vec{e}_1\cdot\Vec{e}_2=$
    %\item$a\Vec{e}_1\cdot \Vec{e}_3=$
    \item$\Vec{e}_2\cdot a\Vec{e}_2=$
    \item$\Vec{e}_3\times \Vec{e}_1=$
    %\item$\|a\Vec{e}_3\times b\Vec{e}_3\|=$
    \item$\|a\Vec{e}_1\times b\Vec{e}_2\|=$
    %\item$\Vec{e}_1\cdot (\Vec{e}_2+a\Vec{e}_3)=$
    \item$\Vec{e}_1\cdot (\Vec{e}_2+a\Vec{e}_1 )=$
    %\item$(a\Vec{e}_1+d\Vec{e}_3)\cdot (b\Vec{e}_2+c\Vec{e}_3)=$
    \item$\|(a\Vec{e}_1+d\Vec{e}_3)\times (b\Vec{e}_2+c\Vec{e}_3)\|=$
\end{enumerate}
\end{multicols}
% --- SOLUCIÓN ---------------------------------------------------------------
\begin{solucion}{Operaciones con base única (primera parte)}
\begin{multicols}{2} 
\begin{enumerate}
    \item $\vec{e}_1 \cdot \vec{e}_1 = \|\vec{e}_1\|^2 = 1$ (por ser base orto-normal).
    \item $a\vec{e}_1 \cdot b\vec{e}_2 = ab(\vec{e}_1 \cdot \vec{e}_2) = ab(0) = 0$ (vectores perpendiculares).
    \item $\vec{e}_2 \times \vec{e}_3 = \vec{e}_1$ (orden cíclico positivo).
    \item $\vec{e}_3 \times \vec{e}_2 = -(\vec{e}_2 \times \vec{e}_3) = -\vec{e}_1$ (propiedad anticonmutativa).
    \item $\vec{e}_1 \times \vec{e}_1 = \vec{0}$ (el producto cruz de un vector consigo mismo es nulo).
    \item $\| a\vec{e}_2 + b\vec{e}_2 \| = \|(a+b)\vec{e}_2\| = |a+b| \cdot 1 = |a+b|$.
    \item $(a\vec{e}_1 + b\vec{e}_2) \cdot \vec{e}_1 = a(\vec{e}_1 \cdot \vec{e}_1) + b(\vec{e}_2 \cdot \vec{e}_1) = a(1) + b(0) = a$.
    \item $\| a\vec{e}_1 \times b\vec{e}_3 \| = \| ab(\vec{e}_1 \times \vec{e}_3) \| = \| ab(-\vec{e}_2) \| = |ab| \cdot 1 = |ab|$.
    \item $\vec{e}_1 \cdot (\vec{e}_2 \times \vec{e}_3) = \vec{e}_1 \cdot \vec{e}_1 = 1$.
    \item $(\vec{e}_1 \times \vec{e}_2) \cdot \vec{e}_3 = \vec{e}_3 \cdot \vec{e}_3 = 1$.
    \item $\| a\vec{e}_1 - a\vec{e}_1 \| = \| \vec{0} \| = 0$.
    \item $(a\vec{e}_1 + b\vec{e}_3) \cdot (c\vec{e}_2) = ac(\vec{e}_1 \cdot \vec{e}_2) + bc(\vec{e}_3 \cdot \vec{e}_2) = 0 + 0 = 0$.
\end{enumerate}
\end{multicols}
\end{solucion}

% --- SOLUCIÓN ---------------------------------------------------------------
\begin{solucion}{Operaciones con base única (segunda parte)}
\begin{multicols}{2}
            
    \begin{enumerate}
        
    \item$\vec{e}_1\cdot\vec{e}_2=0$
    %\item$a\vec{e}_1\cdot \vec{e}_3=$
    \item$\vec{e}_2\cdot a\vec{e}_2= a$
    \item$\vec{e}_3\times \vec{e}_1=\vec{e}_2$
    %\item$\|a\vec{e}_3\times b\vec{e}_3\|=$
    \item$\|a\vec{e}_1\times b\vec{e}_2\|=ab$
    %\item$\vec{e}_1\cdot (\vec{e}_2+a\vec{e}_3)=$
    \item$\vec{e}_1\cdot (\vec{e}_2+a\vec{e}_1 )=a$
    %\item$(a\vec{e}_1+d\vec{e}_3)\cdot (b\vec{e}_2+c\vec{e}_3)=$
    \item$\|(a\vec{e}_1+d\vec{e}_3)\times (b\vec{e}_2+c\vec{e}_3)\|=\sqrt{(ab)^2+(ac)^2+(db)^2}$
        \end{enumerate}
        \end{multicols}
\end{solucion}
